% !TEX program = lualatex
\documentclass[aspectratio=43,professionalfonts,handout]{beamer}
\usepackage{beamer_template}

\newcommand{\LectureNum}{3}
\newcommand{\LectureTitle}{微分法則と積分法則}
\newcommand{\TextPages}{p.\,53-57}
\newcommand{\CourseName}{応用数学}
\renewcommand{\TermName}{前期}

\begin{document}

\begin{frame}{\CourseName\quad 第\LectureNum 回「\LectureTitle」}
  {\small \TermName\quad \TeacherName\quad ／\quad 教科書 \TextPages}

  \textbf{目標}：微分法則・積分法則を理解し、導関数を含む関数のラプラス変換を求めることができる

\end{frame}

\begin{frame}{原関数の微分法則}
  \begin{block}{}
  \begin{itemize}
    \item $\mathcal{L}[f'(t)] = sF(s) - f(0)$
    \item $\mathcal{L}[f''(t)] = s^{2}F(s) - sf(0) - f'(0)$
    \item 高次微分への一般化
    \item 微分方程式への応用への布石
  \end{itemize}
  \end{block}
\end{frame}

\begin{frame}{像関数の微分法則}
  \begin{block}{}
  \begin{itemize}
    \item $\mathcal{L}[tf(t)] = -F'(s)$
    \item $\mathcal{L}[t^n f(t)] = (-1)^n F^{(n)}(s)$
    \item 例：$\mathcal{L}[te^{\alpha t}] = 1/(s-\alpha)^{2}$
    \item 例：$\mathcal{L}[t \sin \omega t] = 2\omega s/(s^{2}+\omega^{2})^{2}$
  \end{itemize}
  \end{block}
\end{frame}

\begin{frame}{積分法則}
  \begin{block}{}
  \begin{itemize}
    \item 原関数：$\mathcal{L}[\int_0^t f d\tau] = F(s)/s$
    \item 像関数：$\mathcal{L}[f(t)/t] = \int_s^\infty F(\sigma)d\sigma$
    \item 例：$\mathcal{L}[\sin \omega t / t] = \pi/2 - \arctan(s/\omega)$
  \end{itemize}
  \end{block}
\end{frame}

\begin{frame}{演習}
  \begin{exampleblock}{自分で解いてみよう}
  \begin{itemize}
    \item 問13〜16より選題
  \end{itemize}
  \end{exampleblock}
\end{frame}

\begin{frame}{まとめ}
  \begin{itemize}
    \item 微分法則（原・像・高次）
    \item 積分法則（原・像）
    \item 次回：逆ラプラス変換
  \end{itemize}
\end{frame}

\end{document}
